% IFMD: Atomic Language Processing (ALP)
% Author: —
% Date: \today
\documentclass[11pt]{article}

% ---- Packages ----
\usepackage[utf8]{inputenc}
\usepackage[T1]{fontenc}
\usepackage{lmodern}
\usepackage[a4paper,margin=1in]{geometry}
\usepackage{microtype}
\usepackage{amsmath,amssymb,mathtools}
\usepackage{bm}
\usepackage{booktabs}
\usepackage{graphicx}
\usepackage{xcolor}
\usepackage{enumitem}
\usepackage{hyperref}
\hypersetup{colorlinks=true,linkcolor=blue,citecolor=blue,urlcolor=blue}
\usepackage{listings}
\lstdefinestyle{code}{
  basicstyle=\ttfamily\small,
  breaklines=true,
  frame=single,
  numbers=left,
  numberstyle=\tiny, 
  keywordstyle=\color{blue!60!black},
  commentstyle=\color{green!50!black},
  stringstyle=\color{red!60!black},
  showstringspaces=false,
  tabsize=2
}
\lstset{style=code,language=Python}

% ---- Macros ----
\newcommand{\N}{\mathbb{N}}
\newcommand{\Z}{\mathbb{Z}}
\newcommand{\R}{\mathbb{R}}
\newcommand{\PP}{\mathcal{P}}
\newcommand{\ord}{\mathrm{ord}}
\newcommand{\sig}{\boldsymbol{\sigma}}
\newcommand{\E}{\mathbb{E}}
\DeclareMathOperator*{\argmin}{arg\,min}
\DeclareMathOperator*{\argmax}{arg\,max}

\title{Atomic Language Processing (ALP) with IFMD: \\
A Three-Layer Prime--Encoded Formalism with Certified Stability}
\author{Ryan O. Van Gelder \& Tyler Van Osdol}
\date{\today}

\begin{document}
\maketitle

\begin{abstract}
We present a practical and mathematically rigorous formulation of \emph{Atomic Language Processing (ALP)} within the IFMD framework. ALP assigns prime-based codes to linguistic units and composes them with prime-encoded tensor signatures (PETC) while enforcing update stability via Arithmetic Control Engine (ACE) constraints. We formalize a three-layer architecture---graphemic (static), grapheme-in-context (lightly dynamic), and lexical/morphemic (dynamic)---define the algebra, derive stability certificates, and provide reference implementations for ranking, isotonic rank adjustment, ACE-constrained assignment, PETC composition, and evaluation.
\end{abstract}

\tableofcontents

\section{Introduction}
The original ALP idea maps high-frequency linguistic relations to small primes and augments them with multiplicity-like features. Directly updating primes additively breaks primality and algebraic meaning. We resolve this by separating (i) an \emph{order prime} determined by rank, from (ii) a \emph{PETC signature} that carries structure via exponents on disjoint feature primes. Learning occurs by updating association scores and re-ranking, while composition is certified by PETC invariants. Stability and reproducibility are enforced using ACE-style spectral constraints.

\paragraph{Contributions.} (1) A clean algebra for order primes and PETC signatures; (2) a three-layer ALP architecture; (3) stable re-ranking via EMA + isotonic adjustment + ACE-constrained assignment; (4) sparse PETC vectors and monoid composition; (5) an evaluation protocol including churn and algebraic-law metrics.

\section{Preliminaries}
\subsection{Prime sets and disjointness}
Let $\PP_\text{ord}$ and $\PP_\text{petc}$ be disjoint sets of primes used for order codes and PETC features, respectively. All integer materializations use
\begin{equation}
  C_\mathbb{Z} = p_{\text{ord}} \cdot \prod_{i} q_i^{e_i}, \quad p_{\text{ord}}\in\PP_\text{ord},\; q_i\in\PP_\text{petc},\; e_i\in\Z_{\ge 0}.
\end{equation}
In-memory, codes remain tuples $(p_{\text{ord}},\,\sig)$ with $\sig\in\Z^{d}_{\ge 0}$ stored sparsely.

\subsection{PETC composition}
A PETC signature is an exponent vector $\sig\in\Z^{d}_{\ge 0}$ over feature primes $\{q_1,\dots,q_d\}$. Composition (tensor product) is vector addition:
\begin{equation}
  \sig_{x\otimes y} = \sig_x + \sig_y. \label{eq:petc-add}
\end{equation}
This yields a commutative monoid $(\Z^{d}_{\ge 0},+)$ with identity $\mathbf{0}$ and supports compile-time validity checks via linear constraints on exponents.

\subsection{Association scores and ranks}
For a unit $u$ (e.g., a grapheme-in-context key), define a robust association score $R(u)$ (e.g., smoothed PPMI or log-likelihood). Online updates use an exponential moving average (EMA):
\begin{equation}
  R^{t+1}(u) = \beta R^{t}(u) + (1-\beta)\,\widehat{R}^{\,t+1}(u), \quad \beta\in[0,1). \label{eq:ema}
\end{equation}
The order prime $p_{\text{ord}}(u)$ is determined by the rank of $R^t$ and reassigned after each update.

\section{Three-Layer ALP}
\subsection{L1: Graphemic layer (static)}
Units are Unicode grapheme clusters. Order primes are fixed per language by corpus frequency. PETC features include position (onset/nucleus/coda), stress, morph role (prefix/root/suffix), and grapheme-to-phoneme (G2P) cluster.

\subsection{L1.5: Grapheme-in-context (dynamic bridge)}
Keys are tuples $u=(g,\,\text{pos},\,\text{stress},\,\text{g2p\_cluster})$ so that domain drift changes ranks meaningfully. After EMA, we compute an \emph{isotonic adjustment} $R^{t+1}_{\text{iso}}$ that preserves near-monotonicity relative to the previous order before re-ranking and ACE-constrained assignment.

\subsection{L2/L3: Morphemic and lexical layers}
Morphemic units (affixes, stems) are lightly dynamic; lexical relations (e.g., lemma--lemma, role arcs) are fully dynamic. PETC captures morphosyntactic features; order primes are assigned by ranks of task-specific association scores under ACE constraints.

\section{Stability via ACE}
Let $A^t$ be a weighted bipartite matrix (units $\times$ contexts) with entries derived from $R^t$. ACE imposes a spectral perturbation budget $\|\Delta A\|_2\le \varepsilon$, which bounds singular value drift by Weyl:
\begin{equation}
  |\sigma_i(A^{t+1})-\sigma_i(A^{t})| \le \|A^{t+1}-A^{t}\|_2 \le \varepsilon.\label{eq:weyl}
\end{equation}
We operationalize this by (i) isotonic rank projection, (ii) min-cost assignment with a churn cap (\(\le K\) swaps), and (iii) rejecting updates that would exceed $\varepsilon$ based on a power-iteration estimate.

\section{Algebraic validity (PETC)}
Mutually exclusive or constrained features are encoded as linear inequalities on $\sig$. For example, a tense conflict is captured by $e_{\text{past}}+e_{\text{future}}\le 1$. A composition $\sig_x+\sig_y$ is valid iff all constraints are satisfied.

\section{Algorithms}
\subsection{Isotonic Rank Adjustment}
Given scores $R$ and a previous order $\pi^{t}$, we compute $R_{\text{iso}}$ by projecting onto the set of vectors that are \emph{nearly} non-increasing along $\pi^{t}$ using pool-adjacent-violators with a slack parameter. This stabilizes ranks prior to assignment.

\subsection{ACE-Constrained Assignment}
We seek a new permutation $\pi^{t+1}$ minimizing total rank displacement subject to a churn budget $K$:
\begin{equation}
  \min_{\pi\in S_n} \sum_{i=1}^n w_{i}\,|\pi(i)-i| \quad \text{s.t.}\quad |\{i: \pi(i)\ne i\}| \le K. \label{eq:assignment}
\end{equation}
A min-cost flow formulation solves (\ref{eq:assignment}) efficiently.

\section{Reference Implementations}
\subsection{Data structures}
\begin{lstlisting}[language=Python,caption={Sparse PETC vector and code tuple}]
from dataclasses import dataclass
from typing import Dict, Tuple

FeatureId = int

@dataclass(frozen=True)
class PETCVec:
    exp: Dict[FeatureId, int]  # feature_id -> exponent (>=0)

    def add(self, other: "PETCVec") -> "PETCVec":
        out = dict(self.exp)
        for k, v in other.exp.items():
            out[k] = out.get(k, 0) + v
        return PETCVec(out)

@dataclass(frozen=True)
class ALPCode:
    order_prime: int
    petc: PETCVec

    def compose(self, other: "ALPCode") -> "ALPCode":
        return ALPCode(
            order_prime=compose_order(self.order_prime, other.order_prime),
            petc=self.petc.add(other.petc)
        )

# Serialization (UF-encoding). PP_order and PP_petc are disjoint prime tables.
PP_petc: Tuple[int, ...] = (...,)  # registry-ordered feature primes

def to_integer(code: ALPCode) -> int:
    n = code.order_prime
    # Multiply only non-zero exponents; avoid huge ints unless needed.
    for idx, q in enumerate(PP_petc):
        e = code.petc.exp.get(idx, 0)
        if e:
            n *= pow(q, e)
    return n
\end{lstlisting}

\subsection{L1 Graphemic encoder}
\begin{lstlisting}[language=Python,caption={Graphemic ALP (static)}]
class GraphemicALP:
    def __init__(self, order_primes: Dict[str, int], feature_map: Dict[str, int]):
        self.order_primes = order_primes          # e.g., {'e':2,'t':3,...}
        self.feature_map = feature_map            # feature name -> feature_id

    def encode(self, grapheme: str, ctx: dict) -> ALPCode:
        p = self.order_primes[grapheme]
        exp = {}
        exp[self.feature_map['position_'+ctx['position']]] = 1   # onset/nucleus/coda
        exp[self.feature_map['stress_'+str(ctx['stress'])]] = 1  # 0/1/2
        exp[self.feature_map['morph_'+ctx['morph_role']]] = 1    # prefix/root/suffix
        exp[self.feature_map['g2p_'+str(ctx['g2p_cluster'])]] = 1
        return ALPCode(order_prime=p, petc=PETCVec(exp))
\end{lstlisting}

\subsection{L1.5 Contextual ranking (EMA + isotonic)}
\begin{lstlisting}[language=Python,caption={EMA + isotonic rank stabilization}]
import numpy as np

def ema_update(R_prev: Dict[tuple, float], R_hat: Dict[tuple, float], beta=0.9):
    R = dict(R_prev)
    for k, v in R_hat.items():
        R[k] = beta * R.get(k, 0.0) + (1 - beta) * v
    return R

def isotonic_adjust(R: Dict[tuple, float], order_prev: list[tuple], slack: float=0.0):
    # Pool-adjacent-violators on values ordered by order_prev; optional slack
    vals = np.array([R[k] for k in order_prev], dtype=float)
    # Simple PAV without slack for brevity
    n = len(vals); i = 0
    while i < n-1:
        if vals[i] < vals[i+1] - slack:
            j = i
            while j >= 0 and vals[j] < vals[i+1] - slack:
                j -= 1
            block = slice(j+1, i+2)
            m = float(np.mean(vals[block]))
            vals[block] = m
            i = max(j, 0)
        else:
            i += 1
    return {k: v for k, v in zip(order_prev, vals)}
\end{lstlisting}

\subsection{ACE-constrained assignment (churn cap)}
\begin{lstlisting}[language=Python,caption={Min-cost assignment with churn cap (sketch)}]
# Given stabilized scores R_iso and previous order, produce a new order with <=K swaps.
# Here we sketch: keep top T fixed, allow limited swaps in a candidate window.

def constrained_order(R_iso: Dict[tuple, float], order_prev: list[tuple], K: int= int(0.05*1e6)):
    # Rank by R_iso
    order_new = sorted(R_iso.keys(), key=lambda k: R_iso[k], reverse=True)
    # Compute minimal edits under cap K (greedy windowed)
    pos_prev = {k:i for i,k in enumerate(order_prev)}
    swaps = 0; final = []
    used = set()
    for k in order_new:
        if k in used: continue
        if k in pos_prev and abs(len(final) - pos_prev[k]) <= 1 and swaps < K:
            final.append(k); used.add(k)
        else:
            final.append(k); used.add(k); swaps += 1
    return final, swaps
\end{lstlisting}

\subsection{Spectral drift (power iteration)}
\begin{lstlisting}[language=Python,caption={Estimate ||ΔA||_2 by power iteration}]
import scipy.sparse as sp
import numpy as np

def spectral_drift_norm(DeltaA: sp.spmatrix, iters: int = 20):
    n = DeltaA.shape[1]
    x = np.random.randn(n)
    x /= np.linalg.norm(x) + 1e-12
    for _ in range(iters):
        y = DeltaA @ x
        normy = np.linalg.norm(y)
        if normy == 0: return 0.0
        x = (DeltaA.T @ y)
        x_norm = np.linalg.norm(x)
        if x_norm == 0: break
        x /= x_norm
    # Rayleigh quotient estimate
    return float(normy)
\end{lstlisting}

\section{Evaluation Protocol}
For each layer, report task metrics and stability/compositional metrics:
\begin{itemize}[leftmargin=2em]
  \item L1: G2P CER, spelling/ocr AUC; PETC-law violations per 1k compositions.
  \item L1.5: prime-churn per update (\%), spectral-drift@k (top-$k$ singular values).
  \item L2: SIGMORPHON inflection F1; contradiction rate (tense/number).
  \item L3: SRL / relation-extraction F1 deltas with ALP features.
\end{itemize}

\section{Serialization and Registries}
Maintain a versioned registry for feature IDs and prime tables. On the wire, emit:
\begin{verbatim}
{
  "version": 1,
  "lang": "en",
  "unit": "grapheme|morpheme|lemma",
  "text_span": {"start": 123, "end": 127},
  "order_prime": 37,
  "petc": {"2":1, "11":2, ...}  // feature_id -> exponent
}
\end{verbatim}

\section{Theoretical Properties}
\paragraph{Monoid.} With composition defined by \eqref{eq:petc-add} and order-prime multiset accumulation, $(\text{ALP},\circ)$ is a monoid with identity (empty code). Associativity follows from integer addition associativity.

\paragraph{Stability certificate.} If $\|\Delta A\|_2\le\varepsilon$, then by \eqref{eq:weyl} all top-$k$ singular values move by at most $\varepsilon$, bounding the rank churn needed to maintain order within a budget.

\section{Example Prime Tables (English, illustrative)}
\begin{table}[h]
\centering
\small
\begin{tabular}{ll|ll|ll}
\toprule
Letter & $p_\text{ord}$ & Letter & $p_\text{ord}$ & Letter & $p_\text{ord}$ \\
\midrule
 e & 2 & t & 3 & a & 5 \\
 o & 7 & i & 11 & n & 13 \\
 s & 17 & h & 19 & r & 23 \\
 d & 29 & l & 31 & c & 37 \\
 u & 41 & m & 43 & w & 47 \\
 f & 53 & g & 59 & y & 61 \\
 p & 67 & b & 71 & v & 73 \\
 k & 79 & j & 83 & x & 89 \\
 q & 97 & z & 101 &  &  \\
\bottomrule
\end{tabular}
\caption{Illustrative static order primes for English letters (per L1).}
\end{table}

\section{Conclusion}
This report consolidates ALP into a compositional, prime-encoded framework with certified stability. The division into static graphemic, dynamic contextual grapheme, and dynamic morpho-lexical layers preserves mathematical elegance (PETC) and practical learning (ACE). The accompanying code sketches provide a starting point for production implementations.

\section*{Acknowledgments}
IFMD framework components: Arithmetic Control Engine (ACE, Tyler Van Osdol) and Prime-Encoded Tensor Calculus (PETC, Ryan Van Gelder).

\begin{thebibliography}{9}
\bibitem{ACE}
Tyler Van Osdol. \emph{Arithmetic Control Engine (ACE)}. Technical memorandum.

\bibitem{PETC}
Ryan Van Gelder. \emph{Prime-Encoded Tensor Calculus (PETC)}. Technical memorandum.
\end{thebibliography}

\appendix
\section{Notes on Implementation Details}
\paragraph{Isotonic slack.} In practice, use a convex projection or isotonic regression with a temperature parameter that trades off fidelity vs.~order preservation.


\paragraph{Assignment solver.} A full min-cost flow with churn constraint yields better optimality than greedy; the sketch provided is a placeholder for clarity.


\paragraph{G2P clusters.} Initialize with a pronouncing dictionary and allow online clustering (e.g., Dirichlet process) over phoneme posteriors to adapt to domain text.


\section{Minimal Runnable Example (Toy)}\label{app:toy}
\noindent\textbf{Goal.} Demonstrate L1.5 dynamic re-ranking (EMA + constrained assignment) and PETC composition on two toy words (\texttt{cat}, \texttt{city}). Pure Python, no external dependencies.


\subsection*{Code}
\begin{lstlisting}[language=Python]
# Minimal ALP toy demo (Python 3.10+)
# - Dynamically ranks grapheme-in-context keys (L1.5)
# - Encodes PETC vectors (sparse) and composes word codes
# - Prints churn and example codes for 'cat' and 'city'


from dataclasses import dataclass
from typing import Dict, Tuple, List
from math import log


# ------------------ Utilities ------------------


def first_n_primes(n: int) -> List[int]:
primes, x = [], 2
def is_prime(k: int) -> bool:
i = 2
while i * i <= k:
if k % i == 0: return False
i += 1
return True
while len(primes) < n:
if is_prime(x): primes.append(x)
x += 1
return primes


# PETC vector (sparse): feature_id -> exponent
@dataclass(frozen=True)
class PETCVec:
exp: Dict[int, int]
def add(self, other: "PETCVec") -> "PETCVec":
out = dict(self.exp)
for k, v in other.exp.items():
out[k] = out.get(k, 0) + v
return PETCVec(out)


@dataclass(frozen=True)
class ALPCode:
order_prime: int
petc: PETCVec
def compose(self, other: "ALPCode") -> "ALPCode":
\end{lstlisting}

\nocite{*}
\bibliographystyle{unsrt}
\bibliography{references}
\end{document}
